% 14-discussion.tex — Résultats 4/4 : discussion (≈ 6 pages).

\chapter{Discussion}
\label{chap:discussion}

\section{Réponses aux questions de recherche}

\begin{arediger}{2 pages}
Reprendre H1 à H3 une par une : confirmée, nuancée ou infirmée, en citant les
chiffres des chapitres~\ref{chap:resultats-quantitatifs}
et~\ref{chap:resultats-qualitatifs}, et confronter à l'état de l'art
(\cite{dejong1998, mayer2009, adams2008phet}).
\end{arediger}

\section{Apports et limites de la méthode de développement}

\begin{arediger}{2 pages}
Ce que l'IA encadrée a permis (volume, qualité vérifiée, cohérence des
notations) et ce qu'elle n'a pas permis (choix de modèle : l'exemple du moteur
abandonné au profit de la barre libre, chapitre~\ref{chap:jeu}, est venu de
l'analyse physique des résultats de simulation, pas de l'IA seule). Temps
réellement passé, compétences requises chez l'enseignant.
\end{arediger}

\section{Validité et transférabilité}

\begin{arediger}{2 pages}
Biais possibles (effet de nouveauté, taille de l'échantillon, auto-sélection,
enseignant-chercheur), validité des instruments, possibilité de réutiliser le
site et la méthode dans d'autres cours.
\end{arediger}
